
\subsection{Objetivo general}

Comparar el desempeño de la comunicación mediante sockets
TCP y gRPC con Protocol Buffers en un entorno distribuido
de tres nodos, a partir de mediciones de latencia y
rendimiento, y comprobar el orden de los eventos mediante
relojes lógicos de Lamport.

\subsection{Objetivos específicos}

\begin{enumerate}

    \item Configurar tres nodos independientes en
    localhost, identificados como nodo-A, nodo-B y
    nodo-C, y verificar el intercambio de mensajes
    mediante sockets TCP y gRPC.

    \item Registrar veinte intercambios por mecanismo
    con sus respectivas marcas de tiempo de Lamport,
    de manera que sea posible analizar el orden
    de envío y recepción de los mensajes.

    \item Medir la latencia de ida y vuelta (RTT)
    con cargas útiles de 100 y 10240 bytes,
    utilizando cinco corridas de cien envíos
    por condición y descartando una corrida
    previa de calentamiento.

    \item Comparar la latencia media, mediana,
    desviación estándar, percentil 95 y rendimiento
    de ambos mecanismos a partir de los registros
    obtenidos durante las pruebas.

    \item Evaluar la comunicación con tres clientes
    concurrentes y verificar las relaciones de
    causalidad observables en los registros de
    eventos.

\end{enumerate}

\subsection{Hipótesis previa a la experimentación}

Se plantea que la implementación con sockets TCP
presentará una latencia media de ida y vuelta menor
que gRPC en las pruebas realizadas sobre localhost,
tanto para mensajes de 100 bytes como de 10240 bytes.

Esta expectativa se relaciona con las características
de las implementaciones desarrolladas. En TCP, el
equipo construyó directamente el intercambio de
solicitudes y respuestas mediante mensajes JSON,
aprovechando el servicio de transporte fiable y
orientado a conexión que proporciona el protocolo
\cite{eddy2022tcp}. En cambio, la implementación
gRPC utiliza llamadas a procedimientos remotos,
mensajes definidos con Protocol Buffers y HTTP/2.
El estudio de Niswar et al. \cite{niswar2024grpc}
analiza el rendimiento de gRPC frente a otras
alternativas de comunicación entre microservicios,
lo que refuerza la necesidad de evaluar estos
mecanismos bajo condiciones experimentales
definidas.

La hipótesis no supone que TCP sea más rápido en
cualquier entorno. Su evaluación se limitará a las
implementaciones del equipo y a las condiciones
establecidas para esta práctica. Se considerará
respaldada cuando la latencia media de TCP sea
inferior a la de gRPC en ambos tamaños de mensaje.
Las diferencias observadas se examinarán junto
con la dispersión y el percentil 95 para evitar
que la comparación dependa únicamente del promedio.
